\documentclass[10pt,a4paper]{amsart}
\usepackage[margin=19mm]{geometry}
\usepackage{amsmath,amssymb,mathtools}
\usepackage[scr=rsfs,cal=euler]{mathalfa}
\usepackage{xfrac}
\usepackage[unicode,hidelinks,bookmarks=false]{hyperref}
\newcommand{\ie}{i.e.\ }
\newcommand{\cf}{cf.\ }
\newcommand{\resp}{respectively\ }
\newcommand{\Subsection}{\subsection}
\providecommand{\nobreakdash}{\nobreak\hbox{-}}
\setlength{\emergencystretch}{2em}
\multlinegap0pt
\usepackage{bm}
\usepackage{bbm}
\usepackage{dsfont}
\usepackage{xspace}
\usepackage{cancel}
\usepackage{mathtools}
\usepackage{amsthm}
\makeatletter
\@ifundefined{theorem}{\newtheorem{theorem}{Theorem}}{}
\@ifundefined{corollary}{\newtheorem{corollary}[theorem]{Corollary}}{}
\@ifundefined{lemma}{\newtheorem{lemma}[theorem]{Lemma}}{}
\makeatother
\def\vp{\varphi}
\title[Gurevic pressure and equidistribution for amenable extension]{Gurevic pressure and equidistribution for amenable extensions of countable state Markov shifts}
\author{Richard Sharp}
\date{Source: arXiv:2510.15663v1 (CC BY 4.0)}
\begin{document}
\maketitle

\noindent\emph{Source and licence.} Gurevic pressure and equidistribution for amenable extensions of countable state Markov shifts. Version arXiv:2510.15663v1 (CC BY 4.0), distributed under the Creative Commons Attribution 4.0 International licence, as recorded in the arXiv licence metadata for this exact version and shown on the abstract page https://arxiv.org/abs/2510.15663v1. Full rights notice: licenses/arxiv-2510.15663-CC-BY-4.0.txt.\par\smallskip

\noindent\textit{Editorial scope.} Counted unit: the Lemma whose source label is \texttt{lem:inequality\_for\_phi} in the article (source0.tex lines 840--859, source label \texttt{lem:inequality\_for\_phi}), delivered together with the article's own complete original proof of it (source0.tex lines 861--965, one \texttt{proof} environment of 105 lines). No mathematics is written, repaired or re-derived: statement and proof are the article's own text line for line. Required context retained uncounted: lem:eigenvalue\_equation (lines 702--707); cor:gamma\_bounded (lines 801--808). Every definition, display and label of those lines is retained unchanged. Omitted from this package and not redistributed: 1--701, 708--800, 809--839, 860--860, 966--1781. Nothing inside a retained excerpt is omitted or altered: every retained body is delivered exactly as the article's lines are written, so no omission receipt of any kind accompanies this package. Citations: the retained excerpts carry no citation command at all, so no bibliography entry of the article is transcribed and no reference is redistributed. Reference adaptation: none was needed and none was made; every reference inside the delivered unit resolves inside it, so no unresolved reference or citation is produced. Printed slips kept literal: no printed word, symbol, hypothesis or number of the counted statement or of its proof was altered. Layout and class: the article's own class is \texttt{amsart} with packages including hyperref, amsmath, amsthm, amssymb, xfrac, faktor, mathtools, mathrsfs, tikz, framed, upgreek, bbm, bm, graphicx; the wrapper is the lane's pinned \texttt{amsart} editorial wrapper at 10pt on A4, which restates the theorem-like environments the retained text needs and adds the article's own macro definitions that the retained text needs (\texttt{\textbackslash vp}), with extra wrapper packages (bm, bbm, dsfont, xspace, cancel, mathtools). The delivered numbering is the unit's own. Assets: no retained span contains a figure environment, a graphics-inclusion command, a PDF-page inclusion, a table or a TikZ picture; no image file is copied into this package, the package declares no asset dependency and no image file is redistributed with it. Licence: version arXiv:2510.15663v1 of this article is distributed under the Creative Commons Attribution 4.0 International licence, as recorded in the arXiv licence metadata for this exact version and shown on the abstract page https://arxiv.org/abs/2510.15663v1. A full rights notice is delivered at licenses/arxiv-2510.15663-CC-BY-4.0.txt. Characters: every retained excerpt is pure ASCII, so the delivered engine, XeLaTeX with its default Latin Modern text font, reports no missing character.
\begin{lemma}\label{lem:eigenvalue_equation}
For all $z \in \mathcal Z \cap [a]$ and all $n \ge 1$,
\[
(\mathcal L_{\vp}^n\Phi)(z,g) = \rho^n\Phi(z,g).
\]
\end{lemma}

\begin{corollary}\label{cor:gamma_bounded}
For all $(z,\gamma) \in \mathcal Z \times G$, we have
\[
0 < \inf_{g \in G} \frac{\Phi(z, g\gamma^{-1})}{\Phi(o,g)} \le
\sup_{g \in G} \frac{\Phi(z,g\gamma^{-1})}{\Phi(o,g)} < 
\infty.
\]
\end{corollary}

\begin{lemma}\label{lem:inequality_for_phi}
There is a function $\phi : \mathcal Z \times G \to \mathbb R^{>0}$ such that
\begin{enumerate}
\item
for all $g \in G$ and $n \ge 1$,
\[
(\mathcal L_{\vp}^n \phi)(o,g) \le \rho^n \phi(o,g),
\]
\item
for all $z \in \mathcal Z$, the function
\[
\gamma \mapsto \frac{\phi(z,\gamma)}{\phi(z,e)}
\]
is a homomorphism from $G$ to $\mathbb R^{>0}$.
\end{enumerate}
Furthermore, for all $a \in S$ and all $\gamma \in G$, we have
\[
\inf_{z \in \mathcal Z \cap [a]} \phi(z,\gamma) >0.
\]
\end{lemma}

\begin{proof}
The construction of the function $\phi$ uses the amenability of $G$.
Let $M$ be a right-invariant Banach mean on $\ell^\infty(G)$. By Jensen's inequality, if $\alpha$ is convex
and $f \in \ell^\infty(G)$
then
\[
M(\alpha(f)) \ge \alpha(M(f)).
\]
We apply this to the functions 
\[
f_{z,\gamma}(g) = \log \frac{\Phi(z,g\gamma^{-1})}{\Phi(o,g)},
\]
where $(z,\gamma) \in \mathcal Z \times G$.
It follows from Corollary \ref{cor:gamma_bounded} that $f_{z,\gamma} \in \ell^\infty(G)$, for all
$(z,\gamma) \in \mathcal Z \times G$. 
Thus, we have
\[
M\left(g \mapsto \frac{\Phi(z,g\gamma^{-1})}{\Phi(o,g)}\right)
= M\left(g \mapsto \exp \log \frac{\Phi(z,g\gamma^{-1})}{\Phi(o,g)}\right)
\ge \exp M(f_{z,\gamma}).
\]
Now define $\phi : \mathcal Z \times G \to \mathbb R^{>0}$ by
\[
\phi(z,\gamma) = \exp M(f_{z,\gamma}).
\]

We will first prove that the inequality in part (i) holds. We need a little care as $M$ is only finitely additive. Thus, we work with finite subsets of $G$ that exhaust $G$.
Let $\{g_k\}_{k=1}^\infty$ be an enumeration of $G$ and, for any $N \ge 1$, let 
$G_N = \{g_1,\ldots,g_N\}$. Then Lemma \ref{lem:eigenvalue_equation} gives us
\begin{align*}
\rho^n &=M\left(g \mapsto \frac{(\mathcal L_\vp^n \Phi)(o,g)}{\Phi(o,g)}\right)
= 
M
\left(g \mapsto \frac{1}{\Phi(o,g)}\sum_{\sigma^n z = o} e^{\vp^n(z)} \Phi(z,g\psi_n(z)^{-1})\right)
\\
&\ge
M
\left(g \mapsto \frac{1}{\Phi(o,g)}\sum_{\substack{\sigma^n z = o \\ \psi_n(z) \in G_N}} 
e^{\vp^n(z)} \Phi(z,g\psi_n(z)^{-1})\right)
\\
&=
\sum_{\substack{\sigma^n z = o \\ \psi_n(z) \in G_N}} e^{\vp^n(z)}
M\left(g \mapsto \frac{\Phi(z,g\psi_n(z)^{-1})}{\Phi(o,g)}\right)
\\
&\ge
\sum_{\substack{\sigma^n z = o \\ \psi_n(z) \in G_N}} e^{\vp^n(z)}
\exp M\left(g \mapsto \log \frac{\Phi(z,g\psi_n(z)^{-1})}{\Phi(o,g)}\right)
\\
&=
\sum_{\substack{\sigma^n z = o \\ \psi_n(z) \in G_N}} e^{\vp^n(z)}
\phi(z,g\psi_n(z)^{-1}).
\end{align*}
Taking the supremum over $N$ gives the inequality in (i).


Next we show that part (ii) holds.
First note that we have
\begin{align*}
\log \phi(o,\gamma \delta) &=
M\left(g \mapsto \log \frac{\Phi(o,g\delta^{-1}\gamma^{-1})}{\Phi(o,g)}\right)
\\
&= M\left(g \mapsto \log \frac{\Phi(o,g\delta^{-1}\gamma^{-1}}{\Phi(o,g\delta^{-1})}\right)
+
M\left(g \mapsto \log \frac{\Phi(o,g\delta^{-1})}{\Phi(o,g)}\right)
\\
&=
\phi(o,\gamma) + \phi(o,\delta)
\end{align*}
and
\begin{align*}
\log \phi(o,\gamma^{-1}) 
&= 
M\left(g \mapsto \log \frac{\Phi(o,g\gamma)}{\Phi(o,g)}\right)
=
M\left(g \mapsto \log \frac{\Phi(o,g)}{\Phi(o,g\gamma^{-1})}\right)
\\
&= 
M\left(g \mapsto -\log \frac{\Phi(o,g\gamma^{-1})}{\Phi(o,g)}\right)
= -\log \phi(o,\gamma).
\end{align*}
Hence $\gamma \mapsto \phi(o,\gamma)$ is a homomorphism.
Next, we observe that
\begin{align*}
\log \phi(z,\gamma) &= 
M\left(g \mapsto \log \frac{\Phi(z,g\gamma)}{\Phi(o,g)}\right)
\\
&= M\left(g \mapsto \log \frac{\Phi(o,g\gamma)}{\Phi(o,g)}\right)
+
M\left(g \mapsto \log \frac{\Phi(z,g\gamma)}{\Phi(o,g\gamma)}\right)
\\
&= M\left(g \mapsto \log \frac{\Phi(o,g\gamma)}{\Phi(o,g)}\right)
+
M\left(g \mapsto \log \frac{\Phi(z,g)}{\Phi(o,g)}\right)
\\
&= \log \phi(o,\gamma) + \log \phi(z,e).
\end{align*}
Hence
\[
\gamma \mapsto \frac{\phi(z,\gamma)}{\phi(z,e)} = \phi(o,\gamma)
\]
is a homomorphism.

To finish, we note that, by Corollary \ref{cor:gamma_bounded}, $\phi(z,\gamma)>0$ and 
one easily sees that this positive lower bound can be taken uniformly in $z \in \mathcal Z \cap [a]$.
\end{proof}

\end{document}
